\section{Execução passo a passo}
\label{sec:passo}

A Tabela~\ref{tab:passo-fib} mostra a execução completa do Fibonacci, no
formato pedido para a apresentação: para cada instrução, o PC, o texto da
linha 1 do LCD, as seis flags (LEDs LD6 a LD1) e $MEM[255]$ (linha 2 do LCD).
Para facilitar o acompanhamento, foram acrescentados os quatro registradores.
Os valores são os vistos ao final de cada instrução, isto é, após o estado
WRITE; os PCs e registradores estão em decimal.

Alguns pontos para observar na tabela:
\begin{itemize}[nosep]
  \item os passos 1 a 3 inicializam \sig{Ra} = 0 e \sig{Rb} = 1 e mostram o
        primeiro termo; a partir do passo 4, cada volta do laço ocupa 7
        instruções (PC 6 a 14);
  \item as flags só mudam nos passos de \instr{ADD}; nas demais instruções
        elas são mantidas, o que é visível nas colunas Z a SM;
  \item no passo 13 a flag EQUAL acende, pois a soma foi $1 + 1$ ($Rb = Ra$);
  \item no passo 90, $233 + 144 = 377$ gera OVERFLOW; no passo 93, \instr{BCS}
        salta para o endereço 16 e o passo 94 executa \instr{HALT}.
\end{itemize}

{\footnotesize
\setlength{\tabcolsep}{4.5pt}
\renewcommand{\arraystretch}{0.95}
\begin{longtable}{rrlcccccc r rrrr}
\caption{Execução passo a passo de \texttt{fibonacci.asm} (valores após o estado WRITE de cada instrução).}\label{tab:passo-fib}\\
\toprule
Passo & PC & LCD linha 1 & Z & N & OVF & EQ & GT & SM & MEM[255] & Ra & Rb & Rc & Rd \\
\midrule\endfirsthead
\toprule
Passo & PC & LCD linha 1 & Z & N & OVF & EQ & GT & SM & MEM[255] & Ra & Rb & Rc & Rd \\
\midrule\endhead
\bottomrule\endfoot
1 & 0 & \texttt{LDI Ra, 0x00} & 0 & 0 & 0 & 0 & 0 & 0 & 000 & 0 & 0 & 0 & 0 \\
2 & 2 & \texttt{LDI Rb, 0x01} & 0 & 0 & 0 & 0 & 0 & 0 & 000 & 0 & 1 & 0 & 0 \\
3 & 4 & \texttt{ST Ra, 0xFF} & 0 & 0 & 0 & 0 & 0 & 0 & 000 & 0 & 1 & 0 & 0 \\
4 & 6 & \texttt{ST Rb, 0xFF} & 0 & 0 & 0 & 0 & 0 & 0 & 001 & 0 & 1 & 0 & 0 \\
5 & 8 & \texttt{MOV Rc, Rb} & 0 & 0 & 0 & 0 & 0 & 0 & 001 & 0 & 1 & 1 & 0 \\
6 & 9 & \texttt{ADD Rb, Ra} & 0 & 0 & 0 & 0 & 1 & 0 & 001 & 0 & 1 & 1 & 0 \\
7 & 10 & \texttt{MOV Ra, Rc} & 0 & 0 & 0 & 0 & 1 & 0 & 001 & 1 & 1 & 1 & 0 \\
8 & 11 & \texttt{LDI Rd, 0x10} & 0 & 0 & 0 & 0 & 1 & 0 & 001 & 1 & 1 & 1 & 16 \\
9 & 13 & \texttt{BCS Rd} & 0 & 0 & 0 & 0 & 1 & 0 & 001 & 1 & 1 & 1 & 16 \\
10 & 14 & \texttt{JMP 0x06} & 0 & 0 & 0 & 0 & 1 & 0 & 001 & 1 & 1 & 1 & 16 \\
11 & 6 & \texttt{ST Rb, 0xFF} & 0 & 0 & 0 & 0 & 1 & 0 & 001 & 1 & 1 & 1 & 16 \\
12 & 8 & \texttt{MOV Rc, Rb} & 0 & 0 & 0 & 0 & 1 & 0 & 001 & 1 & 1 & 1 & 16 \\
13 & 9 & \texttt{ADD Rb, Ra} & 0 & 0 & 0 & 1 & 0 & 0 & 001 & 1 & 2 & 1 & 16 \\
14 & 10 & \texttt{MOV Ra, Rc} & 0 & 0 & 0 & 1 & 0 & 0 & 001 & 1 & 2 & 1 & 16 \\
15 & 11 & \texttt{LDI Rd, 0x10} & 0 & 0 & 0 & 1 & 0 & 0 & 001 & 1 & 2 & 1 & 16 \\
16 & 13 & \texttt{BCS Rd} & 0 & 0 & 0 & 1 & 0 & 0 & 001 & 1 & 2 & 1 & 16 \\
17 & 14 & \texttt{JMP 0x06} & 0 & 0 & 0 & 1 & 0 & 0 & 001 & 1 & 2 & 1 & 16 \\
18 & 6 & \texttt{ST Rb, 0xFF} & 0 & 0 & 0 & 1 & 0 & 0 & 002 & 1 & 2 & 1 & 16 \\
19 & 8 & \texttt{MOV Rc, Rb} & 0 & 0 & 0 & 1 & 0 & 0 & 002 & 1 & 2 & 2 & 16 \\
20 & 9 & \texttt{ADD Rb, Ra} & 0 & 0 & 0 & 0 & 1 & 0 & 002 & 1 & 3 & 2 & 16 \\
21 & 10 & \texttt{MOV Ra, Rc} & 0 & 0 & 0 & 0 & 1 & 0 & 002 & 2 & 3 & 2 & 16 \\
22 & 11 & \texttt{LDI Rd, 0x10} & 0 & 0 & 0 & 0 & 1 & 0 & 002 & 2 & 3 & 2 & 16 \\
23 & 13 & \texttt{BCS Rd} & 0 & 0 & 0 & 0 & 1 & 0 & 002 & 2 & 3 & 2 & 16 \\
24 & 14 & \texttt{JMP 0x06} & 0 & 0 & 0 & 0 & 1 & 0 & 002 & 2 & 3 & 2 & 16 \\
25 & 6 & \texttt{ST Rb, 0xFF} & 0 & 0 & 0 & 0 & 1 & 0 & 003 & 2 & 3 & 2 & 16 \\
26 & 8 & \texttt{MOV Rc, Rb} & 0 & 0 & 0 & 0 & 1 & 0 & 003 & 2 & 3 & 3 & 16 \\
27 & 9 & \texttt{ADD Rb, Ra} & 0 & 0 & 0 & 0 & 1 & 0 & 003 & 2 & 5 & 3 & 16 \\
28 & 10 & \texttt{MOV Ra, Rc} & 0 & 0 & 0 & 0 & 1 & 0 & 003 & 3 & 5 & 3 & 16 \\
29 & 11 & \texttt{LDI Rd, 0x10} & 0 & 0 & 0 & 0 & 1 & 0 & 003 & 3 & 5 & 3 & 16 \\
30 & 13 & \texttt{BCS Rd} & 0 & 0 & 0 & 0 & 1 & 0 & 003 & 3 & 5 & 3 & 16 \\
31 & 14 & \texttt{JMP 0x06} & 0 & 0 & 0 & 0 & 1 & 0 & 003 & 3 & 5 & 3 & 16 \\
32 & 6 & \texttt{ST Rb, 0xFF} & 0 & 0 & 0 & 0 & 1 & 0 & 005 & 3 & 5 & 3 & 16 \\
33 & 8 & \texttt{MOV Rc, Rb} & 0 & 0 & 0 & 0 & 1 & 0 & 005 & 3 & 5 & 5 & 16 \\
34 & 9 & \texttt{ADD Rb, Ra} & 0 & 0 & 0 & 0 & 1 & 0 & 005 & 3 & 8 & 5 & 16 \\
35 & 10 & \texttt{MOV Ra, Rc} & 0 & 0 & 0 & 0 & 1 & 0 & 005 & 5 & 8 & 5 & 16 \\
36 & 11 & \texttt{LDI Rd, 0x10} & 0 & 0 & 0 & 0 & 1 & 0 & 005 & 5 & 8 & 5 & 16 \\
37 & 13 & \texttt{BCS Rd} & 0 & 0 & 0 & 0 & 1 & 0 & 005 & 5 & 8 & 5 & 16 \\
38 & 14 & \texttt{JMP 0x06} & 0 & 0 & 0 & 0 & 1 & 0 & 005 & 5 & 8 & 5 & 16 \\
39 & 6 & \texttt{ST Rb, 0xFF} & 0 & 0 & 0 & 0 & 1 & 0 & 008 & 5 & 8 & 5 & 16 \\
40 & 8 & \texttt{MOV Rc, Rb} & 0 & 0 & 0 & 0 & 1 & 0 & 008 & 5 & 8 & 8 & 16 \\
41 & 9 & \texttt{ADD Rb, Ra} & 0 & 0 & 0 & 0 & 1 & 0 & 008 & 5 & 13 & 8 & 16 \\
42 & 10 & \texttt{MOV Ra, Rc} & 0 & 0 & 0 & 0 & 1 & 0 & 008 & 8 & 13 & 8 & 16 \\
43 & 11 & \texttt{LDI Rd, 0x10} & 0 & 0 & 0 & 0 & 1 & 0 & 008 & 8 & 13 & 8 & 16 \\
44 & 13 & \texttt{BCS Rd} & 0 & 0 & 0 & 0 & 1 & 0 & 008 & 8 & 13 & 8 & 16 \\
45 & 14 & \texttt{JMP 0x06} & 0 & 0 & 0 & 0 & 1 & 0 & 008 & 8 & 13 & 8 & 16 \\
46 & 6 & \texttt{ST Rb, 0xFF} & 0 & 0 & 0 & 0 & 1 & 0 & 013 & 8 & 13 & 8 & 16 \\
47 & 8 & \texttt{MOV Rc, Rb} & 0 & 0 & 0 & 0 & 1 & 0 & 013 & 8 & 13 & 13 & 16 \\
48 & 9 & \texttt{ADD Rb, Ra} & 0 & 0 & 0 & 0 & 1 & 0 & 013 & 8 & 21 & 13 & 16 \\
49 & 10 & \texttt{MOV Ra, Rc} & 0 & 0 & 0 & 0 & 1 & 0 & 013 & 13 & 21 & 13 & 16 \\
50 & 11 & \texttt{LDI Rd, 0x10} & 0 & 0 & 0 & 0 & 1 & 0 & 013 & 13 & 21 & 13 & 16 \\
51 & 13 & \texttt{BCS Rd} & 0 & 0 & 0 & 0 & 1 & 0 & 013 & 13 & 21 & 13 & 16 \\
52 & 14 & \texttt{JMP 0x06} & 0 & 0 & 0 & 0 & 1 & 0 & 013 & 13 & 21 & 13 & 16 \\
53 & 6 & \texttt{ST Rb, 0xFF} & 0 & 0 & 0 & 0 & 1 & 0 & 021 & 13 & 21 & 13 & 16 \\
54 & 8 & \texttt{MOV Rc, Rb} & 0 & 0 & 0 & 0 & 1 & 0 & 021 & 13 & 21 & 21 & 16 \\
55 & 9 & \texttt{ADD Rb, Ra} & 0 & 0 & 0 & 0 & 1 & 0 & 021 & 13 & 34 & 21 & 16 \\
56 & 10 & \texttt{MOV Ra, Rc} & 0 & 0 & 0 & 0 & 1 & 0 & 021 & 21 & 34 & 21 & 16 \\
57 & 11 & \texttt{LDI Rd, 0x10} & 0 & 0 & 0 & 0 & 1 & 0 & 021 & 21 & 34 & 21 & 16 \\
58 & 13 & \texttt{BCS Rd} & 0 & 0 & 0 & 0 & 1 & 0 & 021 & 21 & 34 & 21 & 16 \\
59 & 14 & \texttt{JMP 0x06} & 0 & 0 & 0 & 0 & 1 & 0 & 021 & 21 & 34 & 21 & 16 \\
60 & 6 & \texttt{ST Rb, 0xFF} & 0 & 0 & 0 & 0 & 1 & 0 & 034 & 21 & 34 & 21 & 16 \\
61 & 8 & \texttt{MOV Rc, Rb} & 0 & 0 & 0 & 0 & 1 & 0 & 034 & 21 & 34 & 34 & 16 \\
62 & 9 & \texttt{ADD Rb, Ra} & 0 & 0 & 0 & 0 & 1 & 0 & 034 & 21 & 55 & 34 & 16 \\
63 & 10 & \texttt{MOV Ra, Rc} & 0 & 0 & 0 & 0 & 1 & 0 & 034 & 34 & 55 & 34 & 16 \\
64 & 11 & \texttt{LDI Rd, 0x10} & 0 & 0 & 0 & 0 & 1 & 0 & 034 & 34 & 55 & 34 & 16 \\
65 & 13 & \texttt{BCS Rd} & 0 & 0 & 0 & 0 & 1 & 0 & 034 & 34 & 55 & 34 & 16 \\
66 & 14 & \texttt{JMP 0x06} & 0 & 0 & 0 & 0 & 1 & 0 & 034 & 34 & 55 & 34 & 16 \\
67 & 6 & \texttt{ST Rb, 0xFF} & 0 & 0 & 0 & 0 & 1 & 0 & 055 & 34 & 55 & 34 & 16 \\
68 & 8 & \texttt{MOV Rc, Rb} & 0 & 0 & 0 & 0 & 1 & 0 & 055 & 34 & 55 & 55 & 16 \\
69 & 9 & \texttt{ADD Rb, Ra} & 0 & 0 & 0 & 0 & 1 & 0 & 055 & 34 & 89 & 55 & 16 \\
70 & 10 & \texttt{MOV Ra, Rc} & 0 & 0 & 0 & 0 & 1 & 0 & 055 & 55 & 89 & 55 & 16 \\
71 & 11 & \texttt{LDI Rd, 0x10} & 0 & 0 & 0 & 0 & 1 & 0 & 055 & 55 & 89 & 55 & 16 \\
72 & 13 & \texttt{BCS Rd} & 0 & 0 & 0 & 0 & 1 & 0 & 055 & 55 & 89 & 55 & 16 \\
73 & 14 & \texttt{JMP 0x06} & 0 & 0 & 0 & 0 & 1 & 0 & 055 & 55 & 89 & 55 & 16 \\
74 & 6 & \texttt{ST Rb, 0xFF} & 0 & 0 & 0 & 0 & 1 & 0 & 089 & 55 & 89 & 55 & 16 \\
75 & 8 & \texttt{MOV Rc, Rb} & 0 & 0 & 0 & 0 & 1 & 0 & 089 & 55 & 89 & 89 & 16 \\
76 & 9 & \texttt{ADD Rb, Ra} & 0 & 1 & 0 & 0 & 1 & 0 & 089 & 55 & 144 & 89 & 16 \\
77 & 10 & \texttt{MOV Ra, Rc} & 0 & 1 & 0 & 0 & 1 & 0 & 089 & 89 & 144 & 89 & 16 \\
78 & 11 & \texttt{LDI Rd, 0x10} & 0 & 1 & 0 & 0 & 1 & 0 & 089 & 89 & 144 & 89 & 16 \\
79 & 13 & \texttt{BCS Rd} & 0 & 1 & 0 & 0 & 1 & 0 & 089 & 89 & 144 & 89 & 16 \\
80 & 14 & \texttt{JMP 0x06} & 0 & 1 & 0 & 0 & 1 & 0 & 089 & 89 & 144 & 89 & 16 \\
81 & 6 & \texttt{ST Rb, 0xFF} & 0 & 1 & 0 & 0 & 1 & 0 & 144 & 89 & 144 & 89 & 16 \\
82 & 8 & \texttt{MOV Rc, Rb} & 0 & 1 & 0 & 0 & 1 & 0 & 144 & 89 & 144 & 144 & 16 \\
83 & 9 & \texttt{ADD Rb, Ra} & 0 & 1 & 0 & 0 & 1 & 0 & 144 & 89 & 233 & 144 & 16 \\
84 & 10 & \texttt{MOV Ra, Rc} & 0 & 1 & 0 & 0 & 1 & 0 & 144 & 144 & 233 & 144 & 16 \\
85 & 11 & \texttt{LDI Rd, 0x10} & 0 & 1 & 0 & 0 & 1 & 0 & 144 & 144 & 233 & 144 & 16 \\
86 & 13 & \texttt{BCS Rd} & 0 & 1 & 0 & 0 & 1 & 0 & 144 & 144 & 233 & 144 & 16 \\
87 & 14 & \texttt{JMP 0x06} & 0 & 1 & 0 & 0 & 1 & 0 & 144 & 144 & 233 & 144 & 16 \\
88 & 6 & \texttt{ST Rb, 0xFF} & 0 & 1 & 0 & 0 & 1 & 0 & 233 & 144 & 233 & 144 & 16 \\
89 & 8 & \texttt{MOV Rc, Rb} & 0 & 1 & 0 & 0 & 1 & 0 & 233 & 144 & 233 & 233 & 16 \\
90 & 9 & \texttt{ADD Rb, Ra} & 0 & 0 & 1 & 0 & 1 & 0 & 233 & 144 & 121 & 233 & 16 \\
91 & 10 & \texttt{MOV Ra, Rc} & 0 & 0 & 1 & 0 & 1 & 0 & 233 & 233 & 121 & 233 & 16 \\
92 & 11 & \texttt{LDI Rd, 0x10} & 0 & 0 & 1 & 0 & 1 & 0 & 233 & 233 & 121 & 233 & 16 \\
93 & 13 & \texttt{BCS Rd} & 0 & 0 & 1 & 0 & 1 & 0 & 233 & 233 & 121 & 233 & 16 \\
94 & 16 & \texttt{HALT} & 0 & 0 & 1 & 0 & 1 & 0 & 233 & 233 & 121 & 233 & 16 \\
\end{longtable}

}
